% ============ References (原书 pp.298-315) ============
% 原书文献分两部分: BOOKS (作者缩写键, 如 [M][K][AM][B1-9]) 与 RESEARCH PAPERS
% (按作者字母序, 每位作者内部编号 [1][2]...). 共 435 条 (BOOKS 55 + PAPERS 380),
% 2026-08-08 从原书全部提取, 键经 400dpi 页面图像核验.
% 内部键规则: 书用原键去空格 (如 M, K, AM, G6, N1); 论文用 作者码+编号
% (如 Kap2=Kaplansky 论文2, Kr3=Krull 论文3, Nort8=Northcott 论文8, Coh3=Cohen 论文3).
% 注: 正文中 [B5]/[B6,7]/[B7] 为 Bourbaki 分卷引用 (总条目为 [B1-9]), 保持纯文本不 \\cite.

\section*{Books}
\begin{thebibliography}{XXX}
\bibitem[AB]{AB} M. Auslander and M. Bridger, \textit{Stable module theory}. Memoirs of the AMS 94 (1969).
\bibitem[AM]{AM} M.F. Atiyah and I.G. Macdonald, \textit{Introduction to commutative algebra}. Addison-Wesley, 1969.
\bibitem[An1]{An1} M. Andr\'e, \textit{M\'ethodes simpliciales en alg\`ebre homologique et alg\`ebre commutative}. LNM 32, Springer, 1967.
\bibitem[An2]{An2} M. Andr\'e, \textit{Homologie des alg\`ebres commutatives}. Springer, 1974.
\bibitem[B1-9]{B1-9} N. Bourbaki, \textit{Alg\`ebre commutative}, Chap. 1--Chap. 9. Hermann, 1961--83.
\bibitem[Con]{Con} \textit{Conference on commutative algebra}, Kansas 1972, LNM 311. Springer, 1973.
\bibitem[Diff]{Diff} R. Berger, R. Kiehl, E. Kunz, and H.-J. Nastold, \textit{Differentialrechnung in der analytischen Geometrie}. LNM 38, Springer, 1967.
\bibitem[EG]{EG} E. G. Evans and P. Griffith, \textit{Syzygies}. LMS Lecture Note Series 106, Cambridge University Press, 1985.
\bibitem[F]{F} R. Fossum, \textit{The divisor class group of a Krull domain}. Springer Ergebnisse 74, 1973.
\bibitem[Ful]{Ful} W. Fulton, \textit{Intersection Theory}. Ergebnisse d. Math. 3. Folge Vol. 2, Springer, 1984.
\bibitem[G1]{G1} A. Grothendieck, \textit{\'El\'ements de g\'eom\'etrie alg\'ebrique}, IV. Publ. Math. IHES 20, 1964.
\bibitem[G2]{G2} A. Grothendieck, \textit{\'El\'ements de g\'eom\'etrie alg\'ebrique}, IV. Publ. Math. IHES 24, 1965.
\bibitem[G3]{G3} A. Grothendieck, \textit{\'El\'ements de g\'eom\'etrie alg\'ebrique}, IV. Publ. Math. IHES 28, 1966.
\bibitem[G4]{G4} A. Grothendieck, \textit{\'El\'ements de g\'eom\'etrie alg\'ebrique}, IV. Publ. Math. IHES 32, 1967.
\bibitem[G5]{G5} A. Grothendieck et al., \textit{Cohomologie locale des faisceaux coh\'erents et th\'eor\`emes de Lefschetz locaux et globaux} (SGA 2). North-Holland, 1968.
\bibitem[G6]{G6} A. Grothendieck, \textit{Local cohomology}. LNM 41, Springer, 1967.
\bibitem[GL]{GL} T. H. Gulliksen and G. Levin, \textit{Homology of local rings}. Queen's papers in pure and appl. math., 20, Queen's Univ. Kingston, 1969.
\bibitem[GS]{GS} S. Greco and P. Salmon, \textit{Topics in m-adic topologies}. Springer Ergebnisse 58, 1971.
\bibitem[H]{H} M. Hochster, \textit{Topics in the homological theory of modules over commutative rings}. Regional conference series 24, AMS 1975.
\bibitem[Ha]{Ha} R. Hartshorne, \textit{Algebraic geometry}. Springer, 1977.
\bibitem[HK]{HK} J. Herzog and E. Kunz, \textit{Der kanonische Modul eines Cohen-Macaulay Rings}. LNM 238, Springer, 1971.
\bibitem[HSV]{HSV} M. Hermann, R. Schmidt and W. Vogel, \textit{Theorie der normalen Flachheit}. Teubner Texte zur Math., Leipzig, 1977.
\bibitem[J]{J} N. Jacobson, \textit{Lectures in abstract algebra, Vol. 3. Theory of fields and Galois theory}. Van Nostrand, 1964.
\bibitem[K]{K} I. Kaplansky, \textit{Commutative rings}. Allyn and Bacon, 1970.
\bibitem[Kunz1]{Kunz1} E. Kunz, \textit{Introduction to commutative algebra and algebraic geometry}. Birkh\"auser, 1985.
\bibitem[Kunz2]{Kunz2} E. Kunz, \textit{K\"ahler Differentials}. Vieweg Advanced Lectures in Math., Vieweg \& Sohn, Braunschweig, 1986.
\bibitem[KPR]{KPR} H. Kurke, G. Pfister and M. Roczen, \textit{Henselsche Ringe und algebraische Geometrie}. VEB Deutscher Verlag der Wissensch., Berlin, 1975.
\bibitem[KPPRM]{KPPRM} H. Kurke, G. Pfister, D. Popescu, M. Roczen and T. Mostowski, \textit{Die Approximationseigenschaft lokaler Ringe}. LNM 634, Springer, 1978.
\bibitem[Kr]{Kr} W. Krull, \textit{Idealtheorie}. Springer Ergebnisse, 1935.
\bibitem[La]{La} T.Y. Lam, \textit{Serre's conjecture}. LNM 635, Springer, 1978.
\bibitem[M]{M} H. Matsumura, \textit{Commutative algebra}. Benjamin, 1970 (second edn. 1980).
\bibitem[McA]{McA} S. McAdam, \textit{Asymptotic prime divisors}. LNM 1023, Springer, 1983.
\bibitem[Mac]{Mac} F.S. Macaulay. \textit{Algebraic theory of modular systems}. Cambridge Tracts 19, Cambridge University Press, 1916.
\bibitem[N1]{N1} M. Nagata, \textit{Local rings}. Interscience, 1962.
\bibitem[N2]{N2} M. Nagata, \textit{Lectures on the fourteenth problem of Hilbert}. Tata Inst., 1965.
\bibitem[N3]{N3} M. Nagata, \textit{Polynomial rings and affine spaces}. Regional conference series 37, AMS 1978.
\bibitem[Nor1]{Nor1} D.G. Northcott, \textit{Ideal theory}. Cambridge Tracts 42, Cambridge University Press, 1952.
\bibitem[Nor2]{Nor2} D.G. Northcott, \textit{An introduction to homological algebra}. Cambridge University Press, 1960.
\bibitem[Nor3]{Nor3} D.G. Northcott, \textit{Lessons on rings, modules and multiplicities}. Cambridge University Press, 1968.
\bibitem[Nor4]{Nor4} D.G. Northcott, \textit{Finite free resolutions}. Cambridge Tracts 71, Cambridge University Press, 1976.
\bibitem[R]{R} M. Raynaud, \textit{Anneaux locaux hens\'eliens}. LNM 169, Springer, 1970.
\bibitem[Rob]{Rob} P. Roberts, \textit{Homological invariants of modules over commutative rings}. Presses de l'Univ. de Montr\'eal, 1980.
\bibitem[Rot]{Rot} J. Rotman, \textit{Notes on homological algebra}. Van Nostrand Math Studies 26, 1968.
\bibitem[S1]{S1} P. Samuel, \textit{Alg\`ebre locale}. M\'em. Sci. Math. 123, Gauthier-Villars, 1953.
\bibitem[S2]{S2} P. Samuel, \textit{Lectures on unique factorisation domains}. Tata Inst., 1964.
\bibitem[S3]{S3} P. Samuel, \textit{M\'ethodes d'alg\`ebre abstraite en g\'eom\'etrie alg\'ebrique}. Springer Ergebnisse, 1955.
\bibitem[Sal]{Sal} J. Sally, \textit{Numbers of generators of ideals in local rings}. Lecture Notes in Pure and Appl. Math. 35, Marcel Dekker, 1978.
\bibitem[Sch]{Sch} P. Schenzel, \textit{Dualisierende Komplexe in der lokalen Algebra und Buchsbaum-Ringe}. LNM 907, Springer 1982.
\bibitem[Sh]{Sh} I.R. Shafarevich, \textit{Basic algebraic geometry}. Springer, 1974.
\bibitem[Se]{Se} J.-P. Serre, \textit{Alg\`ebre locale, multiplicit\'es}. LNM 11, Springer, 1965.
\bibitem[St]{St} R.P. Stanley, \textit{Combinatorics and commutative algebra}. Birkh\"auser, 1983.
\bibitem[Su]{Su} S. Suzuki, \textit{Differentials of commutative rings}. Queen's papers in pure and appl. math., 29, Queen's Univ., Kingston, 1971.
\bibitem[SV]{SV} J. St\"uckrad and W. Vogel, \textit{Buchsbaum Rings and Applications}. VEB Deutscher Verlag d. Wissenschaften, Berlin 1986.
\bibitem[V]{V} W. Vasconcelos, \textit{Divisor theory in module categories}. North-Holland Math. Studies 14, North-Holland, 1974.
\bibitem[ZS]{ZS} O. Zariski and P. Samuel, \textit{Commutative algebra}, Vols. I and II. Van Nostrand 1958, 1960 (new edn. Springer, 1975).
\end{thebibliography}

\section*{Research papers}
\begin{thebibliography}{XXX}
\bibitem[1]{Ab1} S. Abhyankar, On the valuations centered in a local domain. Amer. J. Math. 78 (1956), 321--48.
\bibitem[2]{Ab2} S. Abhyankar, Local rings of high embedding dimension. Amer. J. Math. 89 (1967), 1073--7.
\bibitem[1]{AHE1} S. Abhyankar, W. Heinzer and P. Eakin, On the uniqueness of the coefficient ring in a polynomial ring. J. Alg. 23 (1972), 310--42.
\bibitem[1]{Ak1} Y. Akizuki, Einige Bemerkungen \"uber prim\"are Integrit\"atsbereiche mit Teilerkettensatz. Proc. Phys.-Math. Soc. Japan 17 (1935), 327--36.
\bibitem[2]{Ak2} Y. Akizuki, Teilerkettensatz und Vielfachensatz. Proc. Phys.-Math. Soc. Japan 17 (1935), 337--45.
\bibitem[3]{Ak3} Y. Akizuki, Zur Idealtheorie der einartigen Ringbereiche mit dem Teilerkettensatz. Japan. J. Math. 14 (1938), 85--102; II. Japan. J. Math. 14 (1938), 177--87; III. Japan. J. Math. 15 (1939), 1--11.
\bibitem[1]{And1} M. Andr\'e, Localisation de la lissit\'e formelle. manuscripta math. 13 (1974), 297--307.
\bibitem[1]{Ao1} Y. Aoyama, Some basic results on canonical modules. J. Math. Kyoto Univ. 23 (1983), 85--94.
\bibitem[1]{Au1} M. Auslander, Modules over unramified regular local rings. Illinois J. Math. 5 (1961), 631--47.
\bibitem[2]{Au2} M. Auslander, On the purity of the branch locus. Amer. J. Math. 84 (1962), 116--25.
\bibitem[1]{AB1} M. Auslander and D. Buchsbaum, Homological dimension in local rings. Trans. AMS 85 (1957), 390--405.
\bibitem[2]{AB2} M. Auslander and D. Buchsbaum, Codimension and multiplicity. Ann. Math. 68 (1958), 625--57.
\bibitem[3]{AB3} M. Auslander and D. Buchsbaum, Unique factorization in regular local rings. Proc. Natl. Acad. Sci. USA 45 (1959), 733--4.
\bibitem[1]{Av1} L. Avramov, Flat morphisms of complete intersections. Soviet Math. Dokl. 16 (1975), 1413--17.
\bibitem[2]{Av2} L. Avramov, Complete intersections and symmetric algebras. J. Alg. 73 (1981), 248--63.
\bibitem[3]{Av3} L. Avramov, Homology of local flat extensions and complete intersection defects. Math. Ann. 228 (1977), 27--37.
\bibitem[1]{Ba1} H. Bass, On the ubiquity of Gorenstein rings. Math. Z. 82 (1963), 8--28.
\bibitem[1]{Bo1} M. Boratynski, A note on set-theoretical complete intersection ideals. J. Alg. 54 (1978), 1--5.
\bibitem[2]{Bo2} M. Boratynski, Generating ideals up to radical and system of parameters of graded rings. J. Alg. 78 (1982), 20--4.
\bibitem[1]{BI1} A. Brezuleanu and C. Ionescu, On the localization theorems and completion of P-rings. Rev. Roum. Math. Pures Appl. 29 (1984), 371--80.
\bibitem[1]{BR1} A. Brezuleanu and N. Radu, Geometric regularity and formal smoothness. Rend. Accad. Nat. Lincei 51 (1971), 326--7.
\bibitem[2]{BR2} A. Brezuleanu and N. Radu, On the localisation of formal smoothness. Rev. Roum. Math. Pures Appl. 20 (1975), 189--200.
\bibitem[3]{BR3} A. Brezuleanu and N. Radu, Excellent rings and good separation of the module of differentials. Rev. Roum. Math. Pures Appl. 23 (1978), 1455--70.
\bibitem[1]{BRT1} A. Brezuleanu and C. Rotthaus, Eine Bemerkung \"uber Ringe mit geometrisch normalen formalen Fasern. Archiv d. Math. 39 (1982), 19--27.
\bibitem[1]{Br1} M. Brodmann, Rees rings and form rings of almost complete intersections. Nagoya Math. J. 88 (1982), 1--16.
\bibitem[2]{Br2} M. Brodmann, Asymptotic stability of Ass($M/I^nM$). Proc. AMS 74 (1979), 16--18.
\bibitem[1]{BRo1} M. Brodmann and C. Rotthaus, \"Uber den regul\"aren Ort in ausgezeichneten Ringen. Math. Z. 175 (1980), 81--7.
\bibitem[2]{BRo2} M. Brodmann and C. Rotthaus, Local domains with bad sets of formal prime divisors. J. Alg. 75 (1982), 386--94.
\bibitem[1]{Bru1} W. Bruns, `Jede' endliche freie Aufl\"osung ist freie Aufl\"osung eines von drei Elementen erzeugten Ideals. J. Alg. 39 (1976), 429--39.
\bibitem[2]{Bru2} W. Bruns, Zur Konstruktion basischer Elemente. Math. Z. 172 (1980), 63--75.
\bibitem[3]{Bru3} W. Bruns, The Eisenbud-Evans generalized principal ideal theorem and determinantal ideals. Proc. AMS 83 (1981), 19--24.
\bibitem[1]{BE1} D. Buchsbaum and D. Eisenbud, What makes a complex exact? J. Alg. 25 (1973), 259--68.
\bibitem[2]{BE2} D. Buchsbaum and D. Eisenbud, Some structure theorems for finite free resolutions. Adv. in Math. 12 (1974), 84--139.
\bibitem[3]{BE3} D. Buchsbaum and D. Eisenbud, What annihilates a module? J. Alg. 47 (1977), 231--43.
\bibitem[1]{Bu1} L. Burch, On ideals of finite homological dimension in local rings. Proc. Camb. Phil. Soc. 64 (1968), 941--52.
\bibitem[1]{Ca1} A. Caruth, A short proof of the principal ideal theorem. Quart. J. Math. Oxford 31 (1980), 401.
\bibitem[1]{Ch1} S. Chase, Direct product of modules. Trans. AMS 97 (1960), 457--73.
\bibitem[1]{Chv1} C. Chevalley, On the theory of local rings. Ann. Math. 44 (1943), 690--708.
\bibitem[2]{Chv2} C. Chevalley, Intersection of algebraic and algebroid varieties. Trans. AMS 57 (1945), 1--85.
\bibitem[1]{Chi1} L. Chiantini, On form ideals and Artin-Rees condition. manuscripta math. 36 (1981), 125--45.
\bibitem[1]{Chir1} G. Chiriacescu, On the theory of Japanese rings. Rev. Roum. Math. Pures Appl. 27 (1982), 945--8.
\bibitem[1]{Coh1} I.S. Cohen, On the structure and ideal theory of complete local rings. Trans. AMS 59 (1946), 54--106.
\bibitem[2]{Coh2} I.S. Cohen, Commutative rings with restricted minimum condition. Duke Math. J. 17 (1950), 27--42.
\bibitem[3]{Coh3} I.S. Cohen, Lengths of prime ideal chains. Amer. J. Math. 76 (1954), 654--68.
\bibitem[1]{CS1} I.S. Cohen and A. Seidenberg, Prime ideals and integral dependence. Bull. AMS 52 (1946), 252--61.
\bibitem[1]{CN1} C. Cowsik and M.V. Nori, Affine curves in characteristic p are set-theoretic complete intersection. Inv. math. 45 (1978), 111--14.
\bibitem[1]{Da1} E. Davis, Ideals of the principal class, R-sequences and a certain monoidal transformation. Pacific J. Math. 20 (1967), 197--205.
\bibitem[1]{DEP1} C. DeConcini, D. Eisenbud and C. Procesi, Young diagrams and determinantal varieties. Inv. math. 56 (1980), 129--65.
\bibitem[2]{DEP2} C. DeConcini, D. Eisenbud and C. Procesi, Hodge algebras. Ast\'erisque 91 (1982), 1--87.
\bibitem[1]{DP1} C. DeConcini and C. Procesi, A characteristic-free approach to invariant theory. Adv. in Math. 21 (1976), 330--54.
\bibitem[1]{Die1} J. Dieudonn\'e, On regular sequences. Nagoya Math. J. 27 (1966), 355--6.
\bibitem[1]{Du1} S. Dutta, Symbolic powers, intersection multiplicity, and asymptotic behaviour of Tor. J. London Math. Soc. 28 (1983), 261--81.
\bibitem[2]{Du2} S. Dutta, Frobenius and multiplicities. J. Alg. 85 (1983), 424--48.
\bibitem[1]{DHM1} S. Dutta, M. Hochster and J.E. McLaughlin, Modules of finite projective dimension with negative intersection multiplicities. Inv. math. 79 (1985), 253--91.
\bibitem[1]{Ea1} J.A. Eagon, Examples of Cohen-Macaulay rings which are not Gorenstein. Math. Z. 109 (1969), 109--11.
\bibitem[1]{EH1} J.A. Eagon and M. Hochster, R-sequences and indeterminates. Quart. J. Math. Oxford 25 (1974), 61--71.
\bibitem[1]{EN1} J.A. Eagon and D.G. Northcott, Ideals defined by matrices and a certain complex associated with them. Proc. Roy. Soc. London 269 (1962), 188--204.
\bibitem[2]{EN2} J.A. Eagon and D.G. Northcott, Generically acyclic complexes and generically perfect ideals. Proc. Roy. Soc. London 299 (1967), 147--72.
\bibitem[1]{Eak1} P. Eakin, The converse to a well known theorem on Noetherian rings. Math. Ann. 177 (1968), 278--82.
\bibitem[1]{ES1} B. Eckmann and A. Schopf, \"Uber injektiven Moduln. Archiv d. Math. 4 (1953), 75--8.
\bibitem[1]{Ei1} D. Eisenbud, Some directions of recent progress in commutative algebra. Proc. Symposia in Pure Math. 29 (1975), 111--28.
\bibitem[2]{Ei2} D. Eisenbud, Introduction to algebras with straightening laws. In Ring Theory and Algebra III, pp. 243--68. Dekker, 1980.
\bibitem[1]{EB1} D. Eisenbud and D. Buchsbaum, Algebra structures for finite free resolutions, and some structure theorems for ideals of codimension 3. Amer. J. Math. 99 (1977), 447--85.
\bibitem[2]{EB2} D. Eisenbud and D. Buchsbaum, Gorenstein ideals of height 3. In Seminar Eisenbud/Singh/Vogel, Teubner Texte zur Math. Bd. 48, 30--48. Teubner, Leipzig, 1980.
\bibitem[1]{EE1} D. Eisenbud and G. Evans, Every algebraic set in n-space is the intersection of n hypersurfaces. Inv. math. 19 (1973), 107--12.
\bibitem[2]{EE2} D. Eisenbud and G. Evans, Generating modules efficiently: Theorems from algebraic K-theory. J. Alg. 27 (1973), 278--305.
\bibitem[3]{EE3} D. Eisenbud and G. Evans, A generalized principal ideal theorem. Nagoya Math. J. 62 (1976), 41--53.
\bibitem[1]{EG1} G. Evans and P. Griffith, The syzygy problem. Ann. Math. 114 (1981), 323--33.
\bibitem[2]{EG2} G. Evans and P. Griffith, Order ideals of minimal generators. Proc. AMS 86 (1982), 375--8.
\bibitem[1]{Fa1} G. Faltings, Ein einfacher Beweis, da\ss{} geometrische Regularit\"at formale Glattheit impliziert. Arch. d. Math. 30 (1978), 284--5.
\bibitem[2]{Fa2} G. Faltings, \"Uber Annulatoren lokaler Kohomologiegruppen. Arch. d. Math. 30 (1978), 473--6.
\bibitem[3]{Fa3} G. Faltings, Zur Existenz dualisierender Komplexe. Math. Z. 162 (1978), 75--86.
\bibitem[4]{Fa4} G. Faltings, \"Uber Macaulaytizierung. Math. Ann. 238 (1978), 175--92.
\bibitem[1]{Fe1} D. Ferrand, Suite r\'eguli\`ere et intersection compl\`ete. C.R. Acad. Paris 264 (1967), 427--8.
\bibitem[2]{Fe2} D. Ferrand, Trivialisation des modules projectifs. La m\'ethode de Kronecker. J. Pure Appl. Alg. 24 (1982), 261--4.
\bibitem[1]{FR1} D. Ferrand and M. Raynaud, Fibres formelles d'un anneau local noeth\'erien. Ann. Sci. \'Ec. Norm. Sup. 3 (1970), 293--311.
\bibitem[1]{Fi1} M. Fiorentini, On relative regular sequences. J. Alg. 18 (1971), 384--9.
\bibitem[1]{Fit1} H. Fitting, Die Determinantenideale eines Moduls. Jahresberichte Deut. Math. Verein. 46 (1936), 195--228.
\bibitem[1]{Fl1} H. Flenner, Die S\"atze von Bertini f\"ur lokale Ringe. Math. Ann. 229 (1977), 97--111.
\bibitem[1]{Fo1} E. Formanek, Faithful Noetherian modules. Proc. AMS 41 (1973), 381--3.
\bibitem[1]{For1} O. Forster, \"Uber die Anzahl der Erzeugenden eines Ideals in einem Noetherschen Ring. Math. Z. 84 (1964), 80--7.
\bibitem[1]{Fos1} R. Fossum, Duality over Gorenstein rings. Math. Scand. 26 (1970), 165--76.
\bibitem[2]{Fos2} R. Fossum, The structure of indecomposable injective modules. Math. Scand. 36 (1975), 291--312.
\bibitem[3]{Fos3} R. Fossum, Commutative extensions by canonical modules are Gorenstein rings. Proc. AMS 40 (1973), 395--400.
\bibitem[1]{FF1} R. Fossum and H.-B. Foxby, The category of graded modules. Math. Scand. 35 (1974), 288--300.
\bibitem[1]{FFGR1} R. Fossum, H.-B. Foxby, P. Griffith and I. Reiten, Minimal injective resolutions with applications to dualizing modules and Gorenstein modules. Publ. Math. IHES 45 (1976), 193--215.
\bibitem[1]{Fox1} H.-B. Foxby, On the $\mu_i$ in a minimal injective resolution. Math. Scand. 29 (1971), 175--186; II: Math. Scand. 41 (1977), 19--44.
\bibitem[2]{Fox2} H.-B. Foxby, Injective modules under flat base change. Proc. AMS 50 (1975), 23--7.
\bibitem[3]{Fox3} H.-B. Foxby, Isomorphisms between complexes with applications to the homological theory of modules. Math. Scand. 40 (1977), 5--19.
\bibitem[4]{Fox4} H.-B. Foxby, Bounded complexes of flat modules. J. Pure Appl. Alg. 15 (1979), 149--72.
\bibitem[5]{Fox5} H.-B. Foxby, A homological theory of complexes of modules. Preprint series No. 19a, 19b., Copenhagen Univ. Math. Inst., 1981.
\bibitem[1]{FT1} H.-B. Foxby and A. Thorup, Minimal injective resolutions under flat base change. Proc. AMS 67 (1977), 27--31.
\bibitem[1]{Go1} O. Goldman, Hilbert rings and the Hilbert Nullstellensatz. Math. Z. 54 (1951), 136--40.
\bibitem[1]{Got1} S. Goto, On Buchsbaum rings. J. Alg. 67 (1980), 272--9.
\bibitem[2]{Got2} S. Goto, On Buchsbaum rings obtained by gluing. Nagoya Math. J. 83 (1981), 123--35.
\bibitem[3]{Got3} S. Goto, Buchsbaum rings of multiplicity 2. J. Alg. 74 (1982), 494--508.
\bibitem[4]{Got4} S. Goto, On the associated graded rings of parameter ideals in Buchsbaum rings. J. Alg. 85 (1983), 490--534.
\bibitem[5]{Got5} S. Goto, A note on quasi-Buchsbaum rings. Proc. AMS 90 (1984), 511--16.
\bibitem[1]{GS1} S. Goto and Y. Shimoda, On the Rees algebra of Cohen-Macaulay local rings. In R. Draper (ed.), Commutative algebra, analytical methods, pp. 201--231. Dekker, 1982.
\bibitem[1]{GSu1} S. Goto and N. Suzuki, Index of reducibility of parameter ideals in a local ring. J. Alg. 87 (1984), 53--88.
\bibitem[1]{GSW1} S. Goto, N. Suzuki and K. Watanabe, On affine semigroup rings. Japan. J. Math. 2 (1976), 1--12.
\bibitem[1]{GW1} S. Goto and K. Watanabe, On graded rings, I, II. I J. Math. Soc. Japan 30 (1978), 179--213; II (P-graded rings) Tokyo J. Math. 1 (1978), 237--61.
\bibitem[2]{GW2} S. Goto and K. Watanabe, The structure of one-dimensional F-pure rings. J. Alg. 49 (1977), 415--31.
\bibitem[1]{Gr1} S. Greco, Henselization of a ring with respect to an ideal. Trans. AMS 144 (1969), 43--65.
\bibitem[2]{Gr2} S. Greco, Two theorems on excellent rings. Nagoya Math. J. 60 (1976), 139--46.
\bibitem[3]{Gr3} S. Greco, A note on universally catenary rings. Nagoya Math. J. 87 (1982), 95--100.
\bibitem[1]{GM1} S. Greco and M. Marinari, Nagata's criterion and openness of loci for Gorenstein and complete intersections. Math. Z. 160 (1978), 207--16.
\bibitem[1]{GrS1} S. Greco and N. Sankaran, On the separable and algebraic closedness of a Hensel couple in its completion. J. Alg. 39 (1976), 335--48.
\bibitem[1]{GT1} S. Greco and C. Traverso, On seminormal schemes. Compositio Math. 40 (1980), 325--65.
\bibitem[1]{GR1} L. Gruson and M. Raynaud, Crit\`eres de platitude et de projectivit\'e. Inv. math. 13 (1971), 1--89.
\bibitem[1]{Har1} R. Hartshorne, Complete intersections and connectedness. Amer. J. Math. 84 (1962), 497--508.
\bibitem[2]{Har2} R. Hartshorne, A property of A-sequences. Bull. Soc. Math. France 94 (1966), 61--6.
\bibitem[1]{Hap1} L. Harper, On differentiably simple algebras. Trans. AMS 100 (1961), 63--72.
\bibitem[1]{HS1} H. Hasse and F.K. Schmidt, Noch eine Begr\"undung der Theorie der h\"oheren Differential-quotienten in einem algebraischen Funktionenk\"orper einer Unbestimmten. J. reine angew. Math. 177 (1937), 215--37.
\bibitem[1]{He1} W. Heinzer, Some remarks on complete integral closure. J. Australian Math. 9 (1969), 310--14.
\bibitem[2]{He2} W. Heinzer, Quotient overrings of an integral domain. Mathematika 17 (1970), 139--48.
\bibitem[1]{HO1} W. Heinzer and J. Ohm, Locally Noetherian commutative rings. Trans. AMS 158 (1971), 273--84.
\bibitem[2]{HO2} W. Heinzer and J. Ohm, Noetherian intersections of integral domains. Trans. AMS 167 (1972), 291--308.
\bibitem[1]{HI1} M. Herrmann and S. Ikeda, Remarks on lifting of Cohen-Macaulay property. Nagoya Math. J. 92 (1983), 121--32.
\bibitem[2]{HI2} M. Herrmann and S. Ikeda, On the Gorenstein property of Rees algebras. manuscripta math. 59 (1987), 471--90.
\bibitem[1]{HOr1} M. Herrmann and U. Orbanz, Faserdimension von Aufblasungen lokaler Ringe und \"Aquimultiplizit\"at. J. Math. Kyoto Univ. 20 (1980), 651--9.
\bibitem[2]{HOr2} M. Herrmann and U. Orbanz, On equimultiplicity. Math. Proc. Camb. Phil. Soc. 91 (1982), 207--13.
\bibitem[3]{HOr3} M. Herrmann and U. Orbanz, Two notes on flatness. manuscripta math. 40 (1982), 109--33.
\bibitem[1]{HSc1} M. Herrmann and R. Schmidt, Zur Transitivit\"at der normalen Flachheit. Inv. math. 28 (1975), 129--36.
\bibitem[1]{BHe1} B. Herzog, On a relation between the Hilbert functions belonging to a local homomorphism. J. London Math. Soc. 25 (1982), 458--66.
\bibitem[1]{JHe1} J. Herzog, Generators and relations of Abelian semigroups and semigroup rings. manuscripta math. 3 (1970), 175--93.
\bibitem[2]{JHe2} J. Herzog, Ein Cohen-Macaulay-Kriterium mit Anwendungen auf den Konormalenmodul und den Differentialmodul. Math. Z. 163 (1978), 149--62.
\bibitem[3]{JHe3} J. Herzog, Deformationen von Cohen-Macaulay Algebren. J. reine angew. Math. 318 (1980), 83--105.
\bibitem[1]{HK1} J. Herzog and E. Kunz, Die Werthalbgruppe eines lokalen Rings der Dimension 1. S.-B. Heidelberger Akad. Wiss. (1971), 27--67.
\bibitem[1]{HSV1} J. Herzog, A. Simis and W. Vasconcelos, Approximation complexes of blowing-up rings. J. Alg. 74 (1982), 466--93; II. J. Alg. 82 (1983), 53--83.
\bibitem[1]{Hib1} T. Hibi, Every affine graded ring has a Hodge algebra structure. Rend. Sem. Mat. Univers. Politecn. Torino, 44-2 (1986), 27--66.
\bibitem[2]{Hib2} T. Hibi, Distributive lattices, affine semigroup rings and algebras with straightening laws. In Commutative Algebra and Combinatorics, Advanced Studies in Pure Math. 11, 1987, pp. 93--109. North-Holland, Amsterdam.
\bibitem[1]{Hil1} D. Hilbert, \"Uber die Theorie der algebraischen Formen. Math. Ann. 36 (1890), 471--534.
\bibitem[2]{Hil2} D. Hilbert, \"Uber die vollen Invariantensysteme. Math. Ann. 42 (1893), 313--73.
\bibitem[1]{Hir1} H. Hironaka, Resolution of singularities of an algebraic variety over a field of characteristic zero. Ann. Math. 79 (1964), 109--326.
\bibitem[1]{Hoc1} M. Hochster, Rings of invariants of tori, Cohen-Macaulay rings generated by monomials, and polytopes. Ann. Math. 96 (1972), 318--37.
\bibitem[2]{Hoc2} M. Hochster, Non-openness of loci in Noetherian rings. Duke Math. J. 40 (1973), 215--19.
\bibitem[3]{Hoc3} M. Hochster, Contracted ideals from integral extensions of regular rings. Nagoya Math. J. 51 (1973), 25--43.
\bibitem[4]{Hoc4} M. Hochster, Constraints on systems of parameters. In Ring Theory, Proc. Oklahoma Conf., pp. 121--61, Dekker, 1974.
\bibitem[5]{Hoc5} M. Hochster, Cohen-Macaulay rings, combinatorics, and simplicial complexes. In Ring Theory II, Proc. Ohio Conf., pp. 171--223. Dekker, 1977.
\bibitem[6]{Hoc6} M. Hochster, The Zariski-Lipman conjecture for homogeneous complete intersections. Proc. AMS 49 (1975), 261--2.
\bibitem[7]{Hoc7} M. Hochster, Cyclic purity versus purity in excellent Noetherian rings. Trans. AMS 231 (1977), 463--88.
\bibitem[8]{Hoc8} M. Hochster, Some applications of the Frobenius in characteristic 0. Bull. AMS 84 (1978), 886--912.
\bibitem[9]{Hoc9} M. Hochster, Canonical elements in local cohomology modules and the direct summand conjecture. J. Alg. 84 (1983), 503--53.
\bibitem[1]{HE1} M. Hochster and J. Eagon, Cohen-Macaulay rings, invariant theory, and generic perfection of determinantal loci. Amer. J. Math. 93 (1971), 1020--58.
\bibitem[1]{HH1} M. Hochster and C. Huneke, Tightly closed ideals. Bull. AMS 18 (1988), 45--8.
\bibitem[1]{HR1} M. Hochster and L.J. Ratliff, Five theorems on Macaulay rings. Pacific J. Math. 44 (1973), 147--72.
\bibitem[1]{HRo1} M. Hochster and J.L. Roberts, Rings of invariants of reductive groups acting on regular rings are Cohen Macaulay. Adv. in Math. 13 (1974), 115--75.
\bibitem[1]{Hop1} C. Hopkins, Rings with minimum condition for left ideals. Ann. Math. 40 (1939), 712--30.
\bibitem[1]{Hu1} C. Huneke, The theory of d-sequences and powers of ideals. Adv. in Math. 46 (1982), 249--99.
\bibitem[2]{Hu2} C. Huneke, On the symmetric and Rees algebra of an ideal generated by a d-sequence. J. Alg. 62 (1980), 268--75.
\bibitem[3]{Hu3} C. Huneke, A remark concerning multiplicities. Proc. AMS 85 (1982), 331--2.
\bibitem[4]{Hu4} C. Huneke, Linkage and the Koszul homology of ideals. Amer. J. Math. 104 (1982), 1043--62.
\bibitem[5]{Hu5} C. Huneke, Numerical invariants of liaison classes. Inv. Math. 75 (1984), 301--25.
\bibitem[1]{HU1} C. Huneke and B. Ulrich, Divisor class groups and deformations. Amer. J. Math. 107 (1985), 1265--303.
\bibitem[2]{HU2} C. Huneke and B. Ulrich, The structure of linkage. Ann. Math. 126 (1987), 277--334.
\bibitem[1]{Ik1} S. Ikeda, The Cohen-Macaulayness of the Rees algebras of local rings. Nagoya Math. J. 89 (1983), 47--63.
\bibitem[2]{Ik2} S. Ikeda, On the Gorensteinness of Rees algebras over local rings. Nagoya Math. J. 102 (1986).
\bibitem[1]{Is1} F. Ischebeck, Eine Dualit\"at zwischen den Funktoren Ext und Tor. J. Alg. 11 (1969), 510--31.
\bibitem[1]{Iv1} B. Iversen, Amplitude inequalities for complexes. Ann. Sci. \'Ec. Norm. Sup. 10 (1977), 547--58.
\bibitem[1]{Jo1} P. Jothilingam, A note on grade. Nagoya Math. J. 59 (1975), 149--52.
\bibitem[1]{KW1} V. Kac and K. Watanabe, Finite linear groups whose ring of invariants is a complete intersection. Bull. AMS 6 (1982), 221--3.
\bibitem[1]{Kap1} I. Kaplansky, Modules over Dedekind rings and valuation rings. Trans. AMS 72 (1952), 327--40.
\bibitem[2]{Kap2} I. Kaplansky, Projective modules. Ann. Math. 68 (1958), 372--7.
\bibitem[1]{Ke1} M. Kersken, Cousinkomplex und Nennersysteme. Math. Z. 182 (1983), 389--402.
\bibitem[1]{KN1} T. Kimura and H. Niitsuma, On Kunz's conjecture. J. Math. Soc. Japan 34 (1982), 371--8.
\bibitem[1]{Ki1} D. Kirby, Coprimary decomposition of Artinian modules. J. London Math. Soc. 6 (1973), 571--6.
\bibitem[1]{Kiy1} K.H. Kiyek, Anwendung von Ideal-Transformationen. manuscripta math. 34 (1981), 327--53.
\bibitem[1]{Kr1} W. Krull, Primidealketten in allgemeinen Ringbereichen. S.-B. Heidelberger Akad. Wiss. 7 (1928).
\bibitem[2]{Kr2} W. Krull, Ein Satz \"uber prim\"are Integrit\"atsbereiche. Math. Ann. 103 (1930), 450--65.
\bibitem[3]{Kr3} W. Krull, Allgemeine Bewertungstheorie. J. reine angew. Math. 167 (1931), 160--96.
\bibitem[4]{Kr4} W. Krull, \"Uber die Zerlegung der Hauptideale in allgemeinen Ringen. Math. Ann. 105 (1931), 1--14.
\bibitem[5]{Kr5} W. Krull, Beitr\"age zur Arithmetik kommutativer Integrit\"atsbereiche. I. Math. Z. 41 (1936), 545--77; II. Math. Z. 41 (1936), 665--79; III. Math. Z. 42 (1937), 745--66; IV. Math. Z. 42 (1937), 767--73; V. Math. Z. 43 (1938), 768--82; VI. Math. Z. 45 (1939), 1--19; VII. Math. Z. 45 (1940), 319--34; VIII. Math. Z. 48 (1942/43), 533--52.
\bibitem[6]{Kr6} W. Krull, Dimensionstheorie in Stellenringen. J. reine angew. Math. 179 (1938), 204--26.
\bibitem[7]{Kr7} W. Krull, Jacobsonsche Ringe, Hilbertscher Nullstellensatz, Dimensionstheorie. Math. Z. 54 (1951), 354--87.
\bibitem[1]{Ku1} E. Kunz, Characterizations of regular local rings of characteristic p. Amer. J. Math. 91 (1969), 772--84.
\bibitem[2]{Ku2} E. Kunz, On Noetherian rings of characteristic p. Amer. J. Math. 98 (1976), 999--1013.
\bibitem[3]{Ku3} E. Kunz, Almost complete intersections are not Gorenstein rings. J. Alg. 28 (1974), 111--15.
\bibitem[1]{KM1} A. Kustin and M. Miller, Deformation and linkage of Gorenstein algebras. Trans. AMS 284 (1984), 501--23.
\bibitem[2]{KM2} A. Kustin and M. Miller, Tight double linkage of Gorenstein algebras. J. Alg. 95 (1985), 384--97.
\bibitem[1]{KMU1} A. Kustin, M. Miller and B. Ulrich, Linkage theory for algebras with pure resolutions. J. Alg. 102 (1986), 199--228.
\bibitem[1]{LR1} K. Langmann and C. Rotthaus, \"Uber den regul\"aren Ort bei japanischen lokalen Ringen. manuscripta math. 27 (1979), 19--30.
\bibitem[1]{Las1} A. Lascoux, Syzygies des vari\'et\'es d\'eterminantales. Adv. in Math. 30 (1978), 202--37.
\bibitem[1]{LL1} T. Larfeldt and C. Lech, Analytic ramifications and flat couples of local rings. Acta math. 146 (1981), 201--8.
\bibitem[1]{Laz1} D. Lazard, Autour de la platitude. Bull. Soc. Math. France 97 (1969), 81--128.
\bibitem[1]{Le1} C. Lech, Inequalities related to certain couples of local rings. Acta math. 112 (1964), 69--89.
\bibitem[1]{LV1} G. Levin and W. Vasconcelos, Homological dimensions and Macaulay rings. Pacific J. Math. 25 (1968), 315--23.
\bibitem[1]{Li1} S. Lichtenbaum, On the vanishing of Tor in regular local rings. Illinois J. Math. 10 (1966), 220--6.
\bibitem[1]{Lip1} J. Lipman, Free derivation modules on algebraic varieties. Amer. J. Math. 87 (1965), 874--98.
\bibitem[2]{Lip2} J. Lipman, On the Jacobian ideal of the module of differentials. Proc. AMS 21 (1969), 422--6.
\bibitem[1]{MacD1} I.G. MacDonald, Secondary representation of modules over a commutative ring. Symposia Mathematica XI (1973), 23--43.
\bibitem[1]{MS1} I.G. MacDonald and R.Y. Sharp, An elementary proof of the non-vanishing of certain local cohomology modules. Quart. J. Math. Oxford 23 (1972), 197--204.
\bibitem[1]{MacR1} R. MacRae, On the homological dimension of certain ideals. Proc. AMS 14 (1963), 746--50.
\bibitem[2]{MacR2} R. MacRae, On an application of the Fitting invariants. J. Alg. 2 (1965), 153--69.
\bibitem[1]{McA1} S. McAdam, Finite coverings by ideals. In Ring Theory, Proc. Oklahoma Conf., pp. 163--71. Dekker, 1974.
\bibitem[2]{McA2} S. McAdam, Going down. Duke J. Math. 39 (1972), 633--6.
\bibitem[3]{McA3} S. McAdam, Saturated chains in Noetherian rings. Indiana Univ. Math. J. 23 (1974), 719--28.
\bibitem[4]{McA4} S. McAdam, Two applications of asymptotic prime divisors. Proc. AMS 84 (1982), 179--80.
\bibitem[1]{ME1} S. McAdam and P. Eakin, The asymptotic Ass. J. Alg. 61 (1979), 71--81.
\bibitem[1]{Mar1} J. Marot, Sur les anneaux universellement japonais. C. R. Acad. Paris, 277 (1973), 1029--31; C.R. Acad. Paris 278 (1974), 1169--72.
\bibitem[1]{MV1} C. Massaza and P. Valabrega, Sul apertura di luoghi in uno schema localmente noetheriano. Boll. UMI 14-A (1977), 564--74.
\bibitem[1]{Mat1} J. Matijevic, Maximal ideal transforms of Noetherian rings. Proc. AMS 54 (1976), 49--52.
\bibitem[1]{MR1} J. Matijevic and P. Roberts, A conjecture of Nagata on graded Cohen-Macaulay rings. J. Math. Kyoto Univ. 14 (1974), 125--8.
\bibitem[1]{Matl1} E. Matlis, Injective modules over Noetherian rings. Pacific J. Math. 8 (1958), 511--28.
\bibitem[2]{Matl2} E. Matlis, Reflexive domains. J. Alg. 8 (1968), 1--33.
\bibitem[3]{Matl3} E. Matlis, The Koszul complex and duality. Comm. Alg. 1 (1974), 87--144.
\bibitem[4]{Matl4} E. Matlis, The higher properties of R-sequences. J. Alg. 50 (1978), 77--112.
\bibitem[1]{Mats1} H. Matsumura, Formal power series rings over polynomial rings, I. In Number Theory, Algebraic Geometry and Commutative Algebra in honor of Y. Akizuki, pp. 511--20. Kinokuniya, Tokyo, 1973.
\bibitem[2]{Mats2} H. Matsumura, Noetherian rings with many derivations. In Contributions to Algebra (dedicated to E. Kolchin), pp. 279--94. Academic Press, New York, 1977.
\bibitem[3]{Mats3} H. Matsumura, Quasi-coefficient rings. Nagoya Math. J. 68 (1977), 123--30.
\bibitem[4]{Mats4} H. Matsumura, Integrable derivations. Nagoya Math. J. 87 (1982), 227--45.
\bibitem[1]{MK1} N. Mohan Kumar, Complete intersections. J. Math. Kyoto Univ. 17 (1977), 533--9.
\bibitem[2]{MK2} N. Mohan Kumar, On two conjectures about polynomial rings. Inv. math. 46 (1978), 225--36.
\bibitem[1]{Mo1} Y. Mori, On the integral closure of an integral domain, I, II. I Mem. Coll. Sci. Univ. Kyoto 27 (1953), 249--56; II Bull. Kyoto Gakugei Univ. 7 (1955), 19--30.
\bibitem[1]{Mur1} P. Murthy, A note on factorial rings. Arch. d. Math. 15 (1964), 418--21.
\bibitem[1]{Nagr1} K.R. Nagarajan, Groups acting on Noetherian rings. Nieuw Archief voor Wiskunde 16 (1968), 25--29.
\bibitem[1]{Nag1} M. Nagata, On the theory of Henselian rings, I, II, III. I Nagoya Math. J. 5 (1953), 45--57; II Nagoya Math. J. 7 (1954), 1--19; III Mem. Coll. Sci. Univ. Kyoto 32 (1959/60), 93--101.
\bibitem[2]{Nag2} M. Nagata, Basic theorems on general commutative rings. Mem. Coll. Sci. Univ. Kyoto 29 (1955), 59--77.
\bibitem[3]{Nag3} M. Nagata, On the derived normal rings of Noetherian integral domains. Mem. Coll. Sci. Univ. Kyoto 29 (1955), 293--303.
\bibitem[4]{Nag4} M. Nagata, The theory of multiplicity in general local rings. In Proc. Intern. Symp. Tokyo-Nikko 1955, Science Council of Japan, 1956, pp. 191--226.
\bibitem[5]{Nag5} M. Nagata, On the chain problem of prime ideals. Nagoya Math. J. 10 (1956), 51--64.
\bibitem[6]{Nag6} M. Nagata, A Jacobian criterion of simple points. Illinois J. Math. 1 (1957), 427--32.
\bibitem[7]{Nag7} M. Nagata, On the purity of branch loci in regular local rings. Illinois J. Math. 3 (1959), 328--33.
\bibitem[8]{Nag8} M. Nagata, On the closedness of singular loci. Publ. Math. IHES 2 (1959), 29--36.
\bibitem[9]{Nag9} M. Nagata, A type of subrings of a Noetherian ring. J. Math. Kyoto Univ. 8 (1968), 465--7.
\bibitem[10]{Nag10} M. Nagata, Flatness of an extension of a commutative ring. J. Math. Kyoto Univ. 9 (1969), 439--48.
\bibitem[1]{Nak1} Y. Nakai, Theory of differentials in commutative rings. J. Math. Soc. Japan 13 (1961), 68--84.
\bibitem[2]{Nak2} Y. Nakai, High order derivations I. Osaka J. Math. 7 (1970), 1--27.
\bibitem[3]{Nak3} Y. Nakai, On a ring with plenty of high order derivations. J. Math. Kyoto Univ. 13 (1972), 159--64.
\bibitem[4]{Nak4} Y. Nakai, On locally finite iterative higher derivations. Osaka J. Math. 15 (1978), 655--62.
\bibitem[1]{NKI1} Y. Nakai, K. Kosaki and Y. Ishibashi, High order derivations II. J. Sci. Hiroshima Univ. Ser. A-I 34 (1970), 17--27.
\bibitem[1]{No1} E. Noether, Idealtheorie in Ringbereichen. Math. Ann. 83 (1921), 24--66.
\bibitem[2]{No2} E. Noether, Abstrakter Aufbau der Idealtheorie in algebraischen Zahl- und Funktionenk\"orpern. Math. Ann. 96 (1926), 26--61.
\bibitem[1]{Nish1} J. Nishimura, Note on Krull domain. J. Math. Kyoto Univ. 15 (1975), 397--400.
\bibitem[2]{Nish2} J. Nishimura, Note on integral closures of a Noetherian integral domain. J. Math. Kyoto Univ. 16 (1976), 117--22.
\bibitem[3]{Nish3} J. Nishimura, On ideal-adic completion of Noetherian rings. J. Math. Kyoto Univ. 21 (1981), 153--69.
\bibitem[1]{Nom1} M. Nomura, Formal power series rings over polynomial rings, II. In Number Theory, Algebraic Geometry and Commutative Algebra in honor of Y. Akizuki, pp. 521--8. Kinokuniya, Tokyo, 1973.
\bibitem[1]{Nort1} D.G. Northcott, On unmixed ideals in regular local rings. Proc. London Math. Soc. 3 (1953), 20--8.
\bibitem[2]{Nort2} D.G. Northcott, A note on classical ideal theory. Proc. Camb. Phil. Soc. 51 (1955), 766--7.
\bibitem[3]{Nort3} D.G. Northcott, On homogeneous ideals. Proc. Glasgow Math. Assoc. 2 (1955), 105--11.
\bibitem[4]{Nort4} D.G. Northcott, On irreducible ideals in local rings. J. London Math. Soc. 32 (1957), 82--8.
\bibitem[5]{Nort5} D.G. Northcott, Syzygies and specializations. Proc. London Math. Soc. 15 (1965), 1--25.
\bibitem[6]{Nort6} D.G. Northcott, Hilbert functions and the Koszul complex. Bull. London Math. Soc. 2 (1970), 69--72.
\bibitem[7]{Nort7} D.G. Northcott, Grade sensitivity and generic perfection. Proc. London Math. Soc. 20 (1970), 597--618.
\bibitem[8]{Nort8} D.G. Northcott, Injective envelopes and inverse polynomials. J. London Math. Soc. 8 (1974), 290--296.
\bibitem[9]{Nort9} D.G. Northcott, Generalized Koszul complexes and Artinian modules. Quart. J. Math. Oxford 23 (1972), 289--297.
\bibitem[1]{NR1} D.G. Northcott and D. Rees, Reductions of ideals in local rings. Proc. Camb. Phil. Soc. 50 (1954), 145--58.
\bibitem[2]{NR2} D.G. Northcott and D. Rees, A note on reductions of ideals with an application to the generalized Hilbert function. Proc. Camb. Phil. Soc. 50 (1954), 353--9.
\bibitem[3]{NR3} D.G. Northcott and D. Rees, Principal systems. Quart. J. Math. Oxford 8 (1957), 119--27.
\bibitem[4]{NR4} D.G. Northcott and D. Rees, Extensions and simplifications of the theory of regular local rings. J. London Math. Soc. 32 (1957), 367--74.
\bibitem[1]{OC1} L. O'Carroll, Generalized fractions, determinantal maps, and top cohomology modules. J. Pure Appl. Alg. 32 (1984), 59--70.
\bibitem[2]{OC2} L. O'Carroll, On the generalized fractions of Sharp and Zakeri. J. London Math. Soc. 28 (1983), 417--27.
\bibitem[1]{Og1} T. Ogoma, Non-catenary pseudo-geometric normal rings. Japan. J. Math. 6 (1980), 147--63.
\bibitem[2]{Og2} T. Ogoma, Cohen-Macaulay factorial domain is not necessarily Gorenstein. Mem. Fac. Sci. Kochi Univ. (Math.) 3 (1982), 65--74.
\bibitem[3]{Og3} T. Ogoma, Some examples of rings with curious formal fibres. Mem. Fac. Sci. Kochi Univ. Math. 1 (1980), 17--22.
\bibitem[4]{Og4} T. Ogoma, Existence of dualizing complexes. J. Math. Kyoto Univ. 24 (1984), 27--48.
\bibitem[1]{Oh1} J. Ohm, Semi-valuations and groups of divisibility. Can. J. Math. 21 (1969), 576--91.
\bibitem[1]{Or1} U. Orbanz, H\"ohere Derivationen und Regularit\"at. J. reine angew. Math. 262/263 (1973), 194--204.
\bibitem[1]{Pe1} C. Peskine, Une g\'en\'eralisation du `main theorem' de Zariski. Bull. Sci. Math. 90 (1966), 119--27.
\bibitem[1]{PS1} C. Peskine and L. Szpiro, Dimension projective finie et cohomologie locale. Publ. Math. IHES 42 (1973), 47--119.
\bibitem[2]{PS2} C. Peskine and L. Szpiro, Liaison des vari\'et\'es alg\'ebriques. I. Inv. math. 26 (1974), 271--302.
\bibitem[3]{PS3} C. Peskine and L. Szpiro, Syzygies et multiplicit\'es. C. R. Acad. Sci. Paris, Ser. A-B. 278 (1974), 1421--4.
\bibitem[1]{Pf1} G. Pfister, Schlechte Henselsche Ringe. Bull. Acad. Polon. Sci. 25 (1977), 1083--8.
\bibitem[1]{PP1} G. Pfister and D. Popescu, Die strenge Approximationseigenschaft lokaler Ringe. Inv. math. 30 (1975), 145--74.
\bibitem[2]{PP2} G. Pfister and D. Popescu, Die Approximationseigenschaft von Primidealen. Bull. Acad. Polon. Sci. 27 (1979), 771--9.
\bibitem[1]{Po1} D. Popescu, General N\'eron desingularization. Nagoya Math. J. 100 (1985), 97--126.
\bibitem[2]{Po2} D. Popescu, General N\'eron desingularization and approximation. Nagoya Math. J. 104 (1986), 85--115.
\bibitem[1]{Qu1} J. Querre, Sur un th\'eor\`eme de Mori-Nagata. C.R. Acad. Paris 285 (1977), 323--4.
\bibitem[2]{Qu2} J. Querre, Sur une propri\'et\'e des anneaux de Krull. Bull. Sci. Math. France, 95 (1971), 341--54.
\bibitem[1]{Qui1} D. Quillen, Projective modules over polynomial rings. Inv. math. 36 (1976), 167--71.
\bibitem[1]{Rab1} S. Rabinowitch, Zum Hilbertschen Nullstellensatz. Math. Ann. 102 (1929), 520.
\bibitem[1]{Ra1} N. Radu, Un crit\`ere diff\'erentiel de lissit\'e formelle. C.R. Acad. Paris 272 (1971), 1166--8.
\bibitem[2]{Ra2} N. Radu, Une caract\'erisation des alg\`ebres noeth\'eriennes r\'eguli\`eres sur un corps de caract\'eristique z\'ero. C.R. Acad. Paris 270 (1970), 851--3.
\bibitem[3]{Ra3} N. Radu, Sur les alg\`ebres dont le module des diff\'erentielles est plat. Rev. Roum. Math. Pures et Appl. 21 (1976), 933--9.
\bibitem[4]{Ra4} N. Radu, Sur un crit\`ere de lissit\'e formelle. Bull. Math. Soc. Sci. Roum. 21 (1977), 133--5.
\bibitem[5]{Ra5} N. Radu, Sur la structure des alg\`ebres locales noeth\'eriennes formellement lisses. Analele Univ. Bucure\c{s}ti, Mathematics, 29 (1980), 81--4.
\bibitem[1]{Rat1} L.J. Ratliff, On quasi-unmixed local domains, the altitude formula, and the chain condition for prime ideals. I Amer. J. Math. 91 (1969), 508--28; II Amer. J. Math. 92 (1970), 99--144.
\bibitem[2]{Rat2} L.J. Ratliff, Characterizations of catenary rings. Amer. J. Math. 93 (1971), 1070--108.
\bibitem[3]{Rat3} L.J. Ratliff, Catenary rings and the altitude formula. Amer. J. Math. 94 (1972), 458--66.
\bibitem[1]{Re1} D. Rees, Valuations associated with a local ring. I Proc. London Math. Soc. 5 (1955), 107--28; II J. London Math. Soc. 31 (1956), 228--35.
\bibitem[2]{Re2} D. Rees, Valuations associated with ideals. I Proc. London Math. Soc. 6 (1956), 161--74; II J. London Math. Soc. 31 (1956), 221--9.
\bibitem[3]{Re3} D. Rees, Two classical theorems of ideal theory. Proc. Camb. Phil. Soc. 52 (1956), 155--7.
\bibitem[4]{Re4} D. Rees, A theorem of homological algebra. Proc. Camb. Phil. Soc. 52 (1956), 605--10.
\bibitem[5]{Re5} D. Rees, The grade of an ideal or module. Proc. Camb. Phil. Soc. 53 (1957), 28--42.
\bibitem[6]{Re6} D. Rees, A note on form rings and ideals. Mathematika 4 (1957), 51--60.
\bibitem[7]{Re7} D. Rees, A note on general ideals. J. London Math. Soc. 32 (1957), 181--6.
\bibitem[8]{Re8} D. Rees, On a problem of Zariski. Illinois J. Math. 2 (1958), 145--59.
\bibitem[9]{Re9} D. Rees, A note on analytically unramified local rings. J. London Math. Soc. 36 (1961), 24--8.
\bibitem[10]{Re10} D. Rees, Rings associated with ideals and analytic spread. Proc. Camb. Phil. Soc. 89 (1981), 423--32.
\bibitem[1]{RS1} D. Rees and R.Y. Sharp, On a theorem of B. Teissier on multiplicities of ideals in local rings. J. London Math. Soc. 18 (1978), 449--63.
\bibitem[1]{Rei1} G.A. Reisner, Cohen-Macaulay quotients of polynomial rings. Adv. in Math. 21 (1976), 30--49.
\bibitem[1]{RM1} G. Restuccia and H. Matsumura, Integrable derivations in rings of unequal characteristic. Nagoya Math. J. 93 (1984), 173--8.
\bibitem[1]{RV1} L. Robbiano and G. Valla, On normal flatness and normal torsion-freeness. J. Alg. 43 (1976), 552--60.
\bibitem[1]{Rob1} P. Roberts, Two applications of dualizing complexes over local rings. Ann. Sci. \'Ec. Norm. Sup. 9 (1976), 103--6.
\bibitem[2]{Rob2} P. Roberts, The vanishing of intersection multiplicities of perfect complexes. Bull. AMS 13 (1985), 127--30.
\bibitem[3]{Rob3} P. Roberts, Le th\'eor\`eme d'intersection. C. R. Acad. Sci. Paris, 304 (1987), 177--80.
\bibitem[1]{Roth1} C. Rotthaus, Nicht ausgezeichnete, universell japanische Ringe. Math. Z. 152 (1977), 107--25.
\bibitem[2]{Roth2} C. Rotthaus, Universell japanische Ringe mit nicht offenem regul\"arem Ort. Nagoya Math. J. 74 (1979), 123--35.
\bibitem[3]{Roth3} C. Rotthaus, Komplettierung semilokaler quasiausgezeichneter Ringe. Nagoya Math. J. 76 (1979), 173--80.
\bibitem[4]{Roth4} C. Rotthaus, Zur Komplettierung ausgezeichneter Ringe. Math. Ann. 253 (1980), 213--26.
\bibitem[5]{Roth5} C. Rotthaus, Potenzreihenerweiterung und formale Fasern in lokalen Ringen mit Approximationseigenschaft. manuscripta math. 42 (1983), 53--65.
\bibitem[1]{Sally1} J. Sally, On the number of generators of powers of an ideal. Proc. AMS 53 (1975), 24--6.
\bibitem[2]{Sally2} J. Sally, Bounds for numbers of generators of Cohen-Macaulay ideals. Pacific J. Math. 63 (1976), 517--20.
\bibitem[3]{Sally3} J. Sally, On the associated graded ring of a local Cohen-Macaulay ring. J. Math. Kyoto Univ. 17 (1977), 19--21.
\bibitem[4]{Sally4} J. Sally, Cohen-Macaulay local rings of maximal embedding dimension. J. Alg. 56 (1979), 168--83.
\bibitem[5]{Sally5} J. Sally, Reductions, local cohomology and Hilbert functions of local rings. In Commutative Algebra, Durham 1981, pp. 231--41; Cambridge University Press 1982.
\bibitem[1]{Sam1} P. Samuel, La notion de multiplicit\'e en alg\`ebre et en g\'eom\'etrie alg\'ebrique. J. math. pure et appl. 30 (1951), 159--274.
\bibitem[2]{Sam2} P. Samuel, On unique factorization domains. Illinois J. Math. 5 (1961), 1--17.
\bibitem[3]{Sam3} P. Samuel, Sur les anneaux factoriels. Bull. Soc. Math. France 89 (1961), 155--73.
\bibitem[1]{SS1} G. Scheja and U. Storch, Differentielle Eigenschaften der Lokalisierungen analytischer Algebren. Math. Ann. 197 (1972), 137--70.
\bibitem[2]{SS2} G. Scheja and U. Storch, \"Uber Spurfunktionen bei vollst\"andigen Durchschnitten. J. reine angew. Math. 278/279 (1975), 174--90.
\bibitem[1]{Sche1} P. Schenzel, Applications of dualizing complexes to Buchsbaum rings. Adv. in Math. 44 (1982), 61--77.
\bibitem[2]{Sche2} P. Schenzel, Multiplizit\"aten in verallgemeinerten Cohen-Macaulay-Moduln. Math. Nachr. 88 (1979), 295--306.
\bibitem[1]{STC1} P. Schenzel, Trung, Ngo Viet and Cuong, Nguyen Tu, Verallgemeinerte Cohen-Macaulay-Moduln. Math. Nachr. 85 (1978), 57--73.
\bibitem[1]{Sd1} A. Seidenberg, Derivations and integral closure. Pacific J. Math. 16 (1966), 167--73.
\bibitem[1]{Ser1} J.-P. Serre, Sur la dimension homologique des anneaux et des modules noeth\'eriens. Proc. Intern. Symp., Tokyo-Nikko, 1955, Science Council of Japan, 1956, 175--89.
\bibitem[1]{Sey1} H. Seydi, Anneaux hens\'eliens et conditions de cha\^ines. I Bull. Soc. Math. France 98 (1970), 9--31; II C.R. Acad. Paris 270 (1970), 696--8; III Bull. Soc. Math. France 98 (1970), 329--36.
\bibitem[2]{Sey2} H. Seydi, Sur deux th\'eor\`emes d'alg\`ebre commutative. C.R. Acad. Paris 271 (1970), 1105--8.
\bibitem[3]{Sey3} H. Seydi, Sur deux th\'eor\`emes d'alg\`ebre commutative. C.R. Acad. Paris 271 (1970), 1169--72.
\bibitem[4]{Sey4} H. Seydi, Un th\'eor\`eme de descente effective universelle et une application. C.R. Acad. Paris 270 (1970), 801--3.
\bibitem[1]{Shar1} R.Y. Sharp, The Cousin complex for a module over a commutative Noetherian ring. Math. Z. 112 (1969), 340--56.
\bibitem[2]{Shar2} R.Y. Sharp, Gorenstein modules. Math. Z. 115 (1970), 117--39.
\bibitem[3]{Shar3} R.Y. Sharp, Local cohomology theory in commutative algebra. Quart. J. Math. Oxford 21 (1970), 425--34.
\bibitem[4]{Shar4} R.Y. Sharp, Dualizing complexes for commutative Noetherian rings. Math. Proc. Camb. Phil. Soc. 78 (1975), 369--86.
\bibitem[5]{Shar5} R.Y. Sharp, Acceptable rings and homomorphic images of Gorenstein rings. J. Alg. 44 (1977), 246--61.
\bibitem[6]{Shar6} R.Y. Sharp, A commutative Noetherian ring which possesses a dualizing complex is acceptable. Math. Proc. Camb. Phil. Soc. 82 (1977), 197--213.
\bibitem[7]{Shar7} R.Y. Sharp, On the attached prime ideals of certain Artinian local cohomology modules. Proc. Edinburgh Math. Soc. 24 (2) (1981), 9--14.
\bibitem[8]{Shar8} R.Y. Sharp, Secondary representations for injective modules over commutative Noetherian rings. Proc. Edinburgh Math. Soc. 20 (2) (1976), 143--51.
\bibitem[1]{SharV1} R.Y. Sharp and P. V\'amos, The dimension of the tensor product of a finite number of field extensions. J. Pure Appl. Alg. 10 (1977), 249--52.
\bibitem[1]{SZ1} R.Y. Sharp and H. Zakeri, Modules of generalized fractions. Mathematika 29 (1982), 32--41.
\bibitem[2]{SZ2} R.Y. Sharp and H. Zakeri, Modules of generalized fractions and balanced big Cohen-Macaulay modules. In R.Y. Sharp (ed.), Commutative Algebra, Durham 1981, pp. 61--82; Cambridge University Press, 1982.
\bibitem[1]{Shi1} Y. Shimoda, A note on Rees algebras of two-dimensional local domains. J. Math. Kyoto Univ. 19 (1979), 327--33.
\bibitem[1]{Si1} B. Singh, On two numerical criteria for normal flatness in codimension one. Math. Nachr. 83 (1978), 241--6.
\bibitem[2]{Si2} B. Singh, Maximally differential prime ideals in a complete local ring. J. Alg. 82 (1983), 331--9.
\bibitem[1]{Stan1} R. Stanley, Relative invariants of finite groups generated by pseudoreflections. J. Alg. 49 (1977), 134--48.
\bibitem[2]{Stan2} R. Stanley, Hilbert functions of graded algebras. Adv. in Math. 28 (1978), 57--83.
\bibitem[1]{Sto1} U. Storch, Bemerkung zu einem Satz von M. Kneser. Arch. d. Math. 23 (1972), 403--4.
\bibitem[1]{Stu1} J. St\"uckrad, \"Uber die kohomologische Charakterisierung von Buchsbaum-Moduln. Math. Nachr. 95 (1980), 265--72.
\bibitem[1]{StV1} J. St\"uckrad and W. Vogel, Toward a theory of Buchsbaum singularities. Amer. J. Math. 100 (1978), 121--46.
\bibitem[1]{SuN1} N. Suzuki, On the generalized local cohomology and its duality. J. Math. Kyoto Univ. 18 (1978), 71--85.
\bibitem[2]{SuN2} N. Suzuki, On the Koszul complex generated by a system of parameters for a Buchsbaum module. Bull. Dept. General Education of Shizuoka College of Pharmacy 8 (1979), 27--35.
\bibitem[3]{SuN3} N. Suzuki, Canonical duality for Buchsbaum modules -- an application of Goto's lemma on Buchsbaum modules. Bull. Dept. General Education of Shizuoka College of Pharmacy 13 (1984), 47--60.
\bibitem[4]{SuN4} N. Suzuki, On a basic theorem for quasi-Buchsbaum modules. Bull. Dept. General Education of Shizuoka College of Pharmacy 11 (1982), 33--40.
\bibitem[1]{SuS1} S. Suzuki, Some results on Hausdorff m-adic modules and m-adic differentials. J. Math. Kyoto Univ. 2 (1963), 157--82.
\bibitem[2]{SuS2} S. Suzuki, On torsion of the module of differentials of a locality which is a complete intersection. J. Math. Kyoto Univ. 4 (1966), 471--5.
\bibitem[3]{SuS3} S. Suzuki, On N\"agler's numbers of discrete valuation rings. J. Math. Kyoto Univ. 11 (1971), 373--5.
\bibitem[1]{Sw1} R. Swan, The number of generators of a module. Math. Z. 102 (1967), 318--22.
\bibitem[2]{Sw2} R. Swan, On seminormality. J. Alg. 67 (1980), 210--29.
\bibitem[1]{Ta1} H. Tanimoto, Some characterisations of smoothness. J. Math. Kyoto Univ. 23 (1983), 695--706.
\bibitem[2]{Ta2} H. Tanimoto, Smoothness of Noetherian rings. Nagoya Math. J. 95 (1984), 163--79.
\bibitem[1]{Tat1} J. Tate, Homology of Noetherian rings and local rings. Illinois J. Math. 1 (1957), 246--61.
\bibitem[1]{Te1} B. Teissier, Cycles \'evanescents, sections planes et conditions de Whitney. Ast\'erisque 7/8 (1973), 285--362.
\bibitem[2]{Te2} B. Teissier, Sur une \'egalit\'e \`a la Minkowski pour les multiplicit\'es. Ann. Math. 106 (1977), 38--44.
\bibitem[3]{Te3} B. Teissier, On a Minkowski-type inequality for multiplicities -- II. In C.P. Ramanujam -- a tribute, pp. 347--61. Tata Inst. Studies in Math. 8, Springer, 1978.
\bibitem[1]{Tr1} Trung, Ngo Viet, On the symbolic powers of determinantal ideals. J. Alg. 58 (1979), 361--9.
\bibitem[2]{Tr2} Trung, Ngo Viet, A class of imperfect prime ideals having the equality of ordinary and symbolic powers. J. Math. Kyoto Univ. 21 (1981), 239--50.
\bibitem[3]{Tr3} Trung, Ngo Viet, Absolutely superficial sequences. Math. Proc. Camb. Phil. Soc. 93 (1983), 35--47.
\bibitem[1]{Ul1} B. Ulrich, Gorenstein rings as specializations of unique factorization domains. J. Alg. 86 (1984), 129--40.
\bibitem[2]{Ul2} B. Ulrich, Gorenstein rings and modules with high number of generators. Math. Z. 188 (1984), 23--32.
\bibitem[3]{Ul3} B. Ulrich, Liaison and deformation. J. Pure Appl. Alg. 39 (1986), 165--75.
\bibitem[4]{Ul4} B. Ulrich, Rings of invariants and linkage of determinantal ideals. Math. Ann. 274 (1986), 1--17.
\bibitem[1]{Va1} P. Valabrega, On the excellent property for power series rings over polynomial rings. J. Math. Kyoto Univ. 15 (1975), 387--95.
\bibitem[2]{Va2} P. Valabrega, A few theorems on completion of excellent rings. Nagoya Math. J. 61 (1976), 127--33.
\bibitem[3]{Va3} P. Valabrega, Formal fibres and openness of loci. J. Math. Kyoto Univ. 18 (1978), 199--208.
\bibitem[1]{Val1} G. Valla, Certain graded algebras are always Cohen-Macaulay. J. Alg. 42 (1976), 537--48.
\bibitem[2]{Val2} G. Valla, On the symmetric and Rees algebras of an ideal. manuscripta math. 30 (1980), 239--55.
\bibitem[1]{VV1} P. Valabrega and G. Valla, Form rings and regular sequences. Nagoya Math. J. 72 (1978), 93--101.
\bibitem[1]{Ve1} U. Vetter, Generische determinantielle Singularit\"aten: homologische Eigenschaften des Derivationenmoduls. manuscripta math. 45 (1984), 161--91.
\bibitem[1]{Vas1} W. Vasconcelos, Ideals defined by R-sequences. J. Alg. 6 (1967), 309--16.
\bibitem[2]{Vas2} W. Vasconcelos, On finitely generated flat modules. Trans. AMS 138 (1969), 505--12.
\bibitem[1]{Vog1} W. Vogel, \"Uber eine Vermutung von D.A. Buchsbaum. J. Alg. 25 (1973), 106--12.
\bibitem[2]{Vog2} W. Vogel, A non-zero-divisor characterization of Buchsbaum modules. Michigan Math. J. 28 (1981), 147--55.
\bibitem[1]{JW1} J. Watanabe, Some remarks on Cohen-Macaulay rings with many zero-divisors. J. Alg. 39 (1976), 1--14.
\bibitem[2]{JW2} J. Watanabe, A note on Gorenstein rings of embedding codimension three. Nagoya Math. J. 50 (1973), 227--32.
\bibitem[3]{JW3} J. Watanabe, m-full ideals. Nagoya Math. J. 106 (1987), 101--11.
\bibitem[4]{JW4} J. Watanabe, The Dilworth number of artinian rings and finite posets with rank function. In Commutative Algebra and Combinatorics, Advanced Studies in Pure Math. 11 (1987), pp. 303--12. North-Holland, Amsterdam.
\bibitem[1]{KWat1} K. Watanabe, Some examples of one dimensional Gorenstein domains. Nagoya Math. J. 49 (1973), 101--9.
\bibitem[2]{KWat2} K. Watanabe, Certain invariant subrings are Gorenstein, I, II. I Osaka J. Math. 11 (1974), 1--8; II Osaka J. Math. 11 (1974), 379--88.
\bibitem[3]{KWat3} K. Watanabe, Invariant subrings which are complete intersections, I. Invariant subrings of finite Abelian groups. Nagoya Math. J. 77 (1980), 89--98.
\bibitem[1]{WITO1} K. Watanabe, T. Ishikawa, S. Tachibana and K. Otsuka, On tensor products of Gorenstein rings. J. Math. Kyoto Univ. 15 (1975), 387--95.
\bibitem[1]{WR1} K. Watanabe and D. Rotillon, Invariant subrings of $C[X,Y,Z]$ which are complete intersections. manuscripta math. 39 (1982), 339--57.
\end{thebibliography}
